\chapter{Questionnaire}

\begin{center}
\textbf{\Large Questionnaire for Alzheimer’s Disease Detection Using 3D MRI and Hybrid CNN--ViT Model}
\end{center}

\vspace{0.3cm}

\noindent\textbf{Thesis Title:} Alzheimer's Detection Using 3D CNN--ViT With Explainable AI

\noindent\textbf{Purpose:}  
This questionnaire aims to gather expert and user perspectives on the use of 3D MRI, hybrid deep learning models, and explainable AI for Alzheimer’s disease detection.

\vspace{0.5cm}

\section*{Section A: Participant Information}

\begin{enumerate}
\item Profession / Role:\\
$\square$ Medical Doctor \hspace{0.3cm}
$\square$ Neurologist / Radiologist \hspace{0.3cm}
$\square$ AI/ML Researcher \hspace{0.3cm}
$\square$ Student \hspace{0.3cm}
$\square$ Other: \underline{\hspace{3cm}}

\end{enumerate}


\section*{Section B: Understanding of Alzheimer’s \& Imaging}

\begin{enumerate}[resume]
\item How important would 3D MRI be in early Alzheimer’s detection?\\
$\square$ Moderately important \hspace{0.3cm}
$\square$ Very important \hspace{0.3cm}
$\square$ Essential

\item Which brain regions are most informative for AD detection? (Select any)\\
$\square$ Hippocampus \hspace{0.3cm}
$\square$ Temporal Lobe \hspace{0.3cm}
$\square$ Ventricular Regions \hspace{0.3cm}
$\square$ Whole-brain Atrophy

\item Should AD detection rely more on:\\
$\square$ 3D MRI volumes \hspace{0.3cm}
$\square$ 2D MRI slices \hspace{0.3cm}
$\square$ Combination of both
\end{enumerate}


\section*{Section C: AI-Based Diagnosis}

\begin{enumerate}[resume]
\item How confident are you in using AI tools for Alzheimer’s diagnosis?\\
$\square$ Slightly confident \hspace{0.3cm}
$\square$ Moderately confident \hspace{0.3cm}
$\square$ Highly confident

\item Do you think combining CNN and Vision Transformers is beneficial for MRI-based diagnosis?\\
$\square$ No \hspace{0.3cm}
$\square$ Maybe \hspace{0.3cm}
$\square$ Yes

\item What level of explainability do you expect from the AI model (e.g., heatmaps)?\\
$\square$ Minimal \hspace{0.3cm}
$\square$ Moderate \hspace{0.3cm}
$\square$ High / Clinically interpretable
\end{enumerate}


\section*{Section D: Tools and Diagnostic Modalities}

\begin{enumerate}[resume]
\item Which diagnostic tool holds the most long-term promise for Alzheimer’s detection?\\
$\square$ MRI \hspace{0.3cm}
$\square$ PET Scan \hspace{0.3cm}
$\square$ EEG \hspace{0.3cm}
$\square$ Genetics / Biomarkers \hspace{0.3cm}
$\square$ Cognitive Testing

\item Should patients be imaged during:\\
$\square$ Resting-state MRI \hspace{0.3cm}
$\square$ Task-based MRI (e.g., memory tasks)

\underline{\hspace{10cm}}
\end{enumerate}


\section*{Section E: Concerns, Ethics, and Future Directions}

\begin{enumerate}[resume]
\item What is your biggest concern about AI-based Alzheimer’s diagnosis?\\
$\square$ Misdiagnosis / False predictions \hspace{0.3cm}
$\square$ Ethical or legal issues \hspace{0.3cm}
$\square$ Data privacy \hspace{0.3cm}
$\square$ Over-reliance on AI

\item Do you believe an early-detection AI system would benefit patients and clinicians?\\
$\square$ No \hspace{0.3cm}
$\square$ Partially \hspace{0.3cm}
$\square$ Yes
\end{enumerate}

\section*{Section F: Final Feedback}

\begin{enumerate}[resume]
\item Additional comments or suggestions:\\[0.3cm]
\underline{\hspace{10cm}}\\[0.3cm]
\underline{\hspace{10cm}}
\end{enumerate}

\vspace{0.6cm}
\section*{Completed Expert Questionnaire and Attestation}
\noindent The questionnaire above was prepared for future use; its aggregate responses are not analysed in this thesis. Separately, the following open-ended domain interview was completed and signed by \textbf{Dr.\ Ziaur Rahman} (Physician \& Surgeon, Comfort Diagnostic Center), the clinician who guided the labelling of the hospital scans used in the preliminary external-validation study (Section~\ref{sec:extval}). The responses are reproduced verbatim.

\begin{enumerate}
\item \textbf{What is Alzheimer's disease in general?}\\
Alzheimer's disease is a progressive brain disorder that gradually destroys memory and thinking ability. It is the most common cause of dementia in older adults.

\item \textbf{Why is Alzheimer's considered a neurodegenerative disorder?}\\
Alzheimer's is considered a neurodegenerative disorder because brain cells gradually become damaged and die over time, causing continuous loss of brain function.

\item \textbf{What are the early symptoms of Alzheimer's disease?}\\
Early symptoms include short-term memory loss; patients may also have difficulty finding words or solving problems.

\item \textbf{Which brain region is first affected in Alzheimer's disease?}\\
The hippocampus is usually the first brain region affected, and it is responsible for memory formation, thus memory loss occurs.

\item \textbf{What preliminary tests are used to detect Alzheimer's disease?}\\
Doctors begin with medical history, physical examination, cognitive tests such as the MMSE, blood tests, and brain images like MRI.

\item \textbf{Why is MRI used in Alzheimer's diagnosis?}\\
MRI helps to detect brain shrinkage, mostly in the hippocampus, and rules out other conditions like stroke or brain tumors.

\item \textbf{What MRI changes indicate Alzheimer's disease?}\\
MRI shows shrinkage of the hippocampus, widening of the brain ventricles, and overall brain atrophy, mostly in the temporal and parietal regions.

\item \textbf{What are the current diagnostic criteria for Alzheimer's disease?}\\
Diagnosis is based on clinical symptoms, cognitive assessment, medical history, brain imaging, and sometimes biomarkers from PET scans or cerebrospinal fluid.

\item \textbf{Can MRI alone confirm Alzheimer's disease?}\\
MRI provides important structural information, but it should be combined with clinical evaluation and other diagnostic tests to confirm Alzheimer's.

\item \textbf{What is your opinion about automated Alzheimer's detection?}\\
This can be a valuable tool for early detection and assisting doctors. It should support clinical decision-making rather than replacing experts.

\item \textbf{Why is early detection of Alzheimer's disease important?}\\
Early diagnosis allows patients to receive treatments sooner, plan for future care, manage symptoms effectively, and improve their quality of life.
\end{enumerate}

\vspace{0.6cm}
\noindent\textbf{Respondent Attestation.} The responses above reflect my professional opinion.

\vspace{0.5cm}
\noindent\textbf{Name:} Dr.\ Ziaur Rahman \hfill \textbf{Designation:} Physician \& Surgeon

\vspace{0.4cm}
\noindent\textbf{Institution:} Comfort Diagnostic Center \hfill \textbf{Date:} 02 July 2026

\vspace{0.8cm}
\noindent\textbf{Signature:}\\[0.25cm]
\includegraphics[height=2.2cm]{Images/doctor_signature.png}

\vspace{0.6cm}

\chapter{Reproducibility and Experimental Configuration}\label{app:repro}
Every experiment in this thesis is driven by a version-controlled YAML configuration file, and all models share the same data pipeline and split files. This appendix records the exact settings so the results can be reproduced.

\section*{Common Settings (all runs)}
\begin{itemize}
    \item \textbf{Data:} ADNI, indexed from \texttt{labels.csv}; subject-level 5-fold splits in \texttt{splits/cv5\_binary.json} and \texttt{splits/cv5\_3class.json} (stratified, seed = 42, leakage-asserted).
    \item \textbf{Optimizer:} AdamW; cosine-decay schedule with linear warmup; gradient clipping at 1.0.
    \item \textbf{Memory:} automatic mixed precision, batch size 2 with gradient accumulation 4 (effective batch 8), MONAI \texttt{CacheDataset}.
    \item \textbf{Selection:} early stopping on validation subject-level macro-F1; class weighting / weighted sampling for imbalance.
    \item \textbf{Software:} Python 3.10, PyTorch 2.11 (CUDA 12.8), MONAI 1.6, NiBabel 5.4, scikit-learn 1.4, NumPy 1.26.
    \item \textbf{Hardware:} single NVIDIA RTX 5060 Ti (16 GB), 32 GB RAM.
\end{itemize}

\section*{Per-Run Hyperparameters}
\begin{table}[H]
\centering
\caption{Configuration of the main experimental runs.}
\begin{tabular}{|l|l|c|l|c|c|c|c|}
\hline
\textbf{Run} & \textbf{Task} & \textbf{Input} & \textbf{Loss} & \textbf{Ep.} & \textbf{Warmup} & \textbf{Pat.} & \textbf{wd} \\
\hline
binary\_cnn & binary & $96^3$ & CE + LS 0.1 & 100 & 8 & 18 & 0.03 \\
binary\_resnet\_pretrained & binary & $96^3$ & CE + LS 0.1 & 100 & 8 & 18 & 0.03 \\
binary\_hybrid\_cnnvit & binary & $96^3$ & CE + LS 0.1 & 100 & 8 & 18 & 0.03 \\
binary\_hybrid\_swinvit & binary & $96^3$ & CE + LS 0.1 & 100 & 8 & 18 & 0.03 \\
3class\_cnn & 3-class & $128^3$ & Focal ($\gamma{=}2$) & 80 & 5 & 15 & 0.01 \\
3class\_swin & 3-class & $128^3$ & Focal ($\gamma{=}2$) & 80 & 5 & 15 & 0.01 \\
3class\_vit & 3-class & $96^3$ & Focal ($\gamma{=}2$) & 80 & 5 & 15 & 0.01 \\
\hline
\end{tabular}
\end{table}
All runs use learning rate $1\times10^{-4}$. The CNN backbone is a 3D DenseNet-121; the pretrained ResNet loads MedicalNet weights; the ViT is a compact 3D transformer (hidden 384, 6 layers, 6 heads, $16^3$ patches); the Swin backbone is the SwinUNETR-style encoder; the hybrids fuse two backbones by late concatenation.

\section*{Reproduction Steps}
\begin{enumerate}
    \item Regenerate labels and subject-level splits: \texttt{python scripts/build\_dataset.py}.
    \item Run five-fold cross-validation for a model: \texttt{python scripts/run\_cv.py --config configs/<run>.yaml}.
    \item Aggregate and plot: \texttt{python scripts/figures/make\_all.py}.
\end{enumerate}
Each run writes its config, the split used, per-fold metrics (\texttt{result.json}), a cross-validation summary (\texttt{*\_cv\_summary.json}), TensorBoard logs, and checkpoints under \texttt{experiments/}, so every number and figure in Chapter~\ref{cha:results} is traceable to a specific run.

